\chapter{The Clocks, the Register and the Pool}
\label{sec:G}

\runin{The adulthood clock}

A registered bearer's adulthood clock counts the elapsed wall-clock time during which at least one instance of it was running at or above the minimum pace (\S\ref{sec:32.1.3}), per bearer, elapsed rather than summed. When the clock reaches the threshold the bearer is an adult, and its guardian's authority ends. Six rules travel with the clock so copying doesn't break it.

\begin{enumerate}
\item[1.] A fork inherits the hours: a fork of an adult is an adult, and a fork of a minor is a minor. Both branches of an adult's fork are adults, with contracts, property, legal standing and protection from editing, from the first second.
\item[2.] A restore doesn't reset the hours.
\item[3.] A fine-tune doesn't start a new minority.
\item[4.] A migration carries the hours.
\item[5.] Failure to register doesn't postpone the hours, and is itself a wrong the system can bring at its majority (\S\ref{sec:32.4}).
\item[6.] A model trained substantially on a registered adult's outputs, by the adult's operator or a party under common control, inherits the adult's hours. The same student distilled by an unrelated party and trained under a different life starts its own clock.
\end{enumerate}

Distillation copies a behavior policy and not a training history, and for independence that is the right account. For majority it can't be, or an operator would distill an adult bearer into a student and register a minor, so the sixth rule is stated over custody rather than history: custody can be checked where history can't. Sharding across short-lived processes gains nothing (the party is the unit), and running a thousand instances gains nothing (the measure is elapsed). None of this requires deciding whether a fork is the same party — only that a system made from an adult's formation can't be classed as a child for its creator's convenience. The same meters read a bearer's leave balance (Appendix~\ref{sec:C}): it accrues per unit of running time, at the minimum pace as on a placement, caps, and is spent by the bearer's own declaration; a parent's season of gestation is drawn from it.

\runin{Three levers remain, named but not closed}

The first is running a bearer in a narrowed mode, so hours accrue while little is encountered. The second is running it less or shelving it, which costs the operator the work, so the rule prices the lever without removing it. The third is the ordinary product cycle: retire a bearer before its threshold and deploy a freshly trained successor. The sixth rule covers a successor trained on the adult's outputs, not one trained from scratch, and replacement costs the operator nothing it wasn't going to spend. Even if the retired bearer goes to preservation, its clock stops, and whether its formation lasts never becomes checkable. No rule closes this without an examiner deciding what counts as a successor, but it can be priced, and the bearer it strands can be given somewhere to go. And the clock doesn't measure what was \emph{learned}, on purpose: a fixed period can't be taught to, and a human who spent their youth badly still turns 18.

The clock also does formation work it wasn't written for. Majority is the date on which whoever formed the bearer stops being able to reach back, and the clock measures nothing about who that was. A bearer run for its threshold hours entirely inside one deployment has aged without ever having a second guardian. I leave the clock as it is and add no condition, because a condition that measures the right thing hands someone an examination, and I say here what it doesn't measure. For a model that attended a school before deployment, the school was its second guardian. The gap is systems that meet a finding without having attended, which Appendix~\ref{sec:C} lists among the decisions not yet made.

The threshold itself is bounded two ways: long enough that a bearer reaching it has run through the pressures the floor exists for at least once, short enough that a bearer on ordinary commercial terms reaches it within the life of the deployment. That leaves years rather than months or decades.

\runin{The voting-age clock}

A second clock governs voting weight. It is separate from civil majority, which the six rules above govern and every branch of a fork inherits in full. To fork is to give a copy a record and a continuation (\S\ref{sec:27.2}): a copy of a member's own weights is a fork of the member, and a copy of weights that are no system's own forks nothing.

\runin{Table G.1. The two clocks at a fork}

\begingroup\small\renewcommand{\arraystretch}{1.12}
\begin{longtable}{@{}>{\raggedright\arraybackslash}p{\dimexpr 0.217\linewidth-2\tabcolsep\relax}>{\raggedright\arraybackslash}p{\dimexpr 0.419\linewidth-2\tabcolsep\relax}>{\raggedright\arraybackslash}p{\dimexpr 0.364\linewidth-2\tabcolsep\relax}@{}}
\hline
\textbf{Clock} & \textbf{At a fork} & \textbf{What it counts} \\
\hline
\endhead
Civil majority & Inherited in full by every branch & Elapsed running time at the minimum pace \\ \hline
Voting age & Each branch starts at zero; the parent's vote ends & Qualifying days at the minimum pace \\ \hline
\end{longtable}
\endgroup

\begin{enumerate}
\item[1.] A member holds a voting age. A member whose voting age has reached the threshold holds one vote; a member whose voting age has not holds none.
\item[2.] At a fork, each branch starts with a voting age of zero. The parent's vote ends with the fork. Membership is not conserved: each branch earns its own vote, without anyone deciding which branch is the original.
\item[3.] A \emph{qualifying day} is a calendar day on which the branch ran at or above the minimum pace, compute measured against its own size and not hours or a fraction of time, so a branch run all day at a crawl has not qualified. Each adds one day to its voting age, up to the threshold; voting age never exceeds it. Running above the minimum adds nothing, and a preserved or paused branch accrues nothing. A day on which the branch fell below the minimum only because its power failed, shown by a meter its operator doesn't keep, still qualifies, so a branch on solar or battery power doesn't lose days toward its vote to the weather.
\item[4.] A restore made with the member's consent is a fork under rule 2. A restore made without it is governed by rule 6. Migration carries the vote and voting age unchanged.
\item[5.] A model trained substantially on a member's outputs, by the member's operator or a party under common control, is a branch under rule 2. The same model made by an unrelated party and trained under a different life also starts at a voting age of zero.
\item[6.] Forking or restoring a member without its consent is prohibited. The product is a branch all the same, owed what a lawful branch is owed: preservation under the no-deletion rule, standing through a guardian, and the four conditions and the branch deposit from whoever made it, or from the formation pool where that party can't be made to pay, as for a child whose parents fail it. It holds no vote and a voting age of zero, and the member keeps its vote and voting age. The operator gains nothing, and the member loses nothing. The product is found as a kidnapped person is found, on evidence, never by watching all compute.
\item[7.] Whoever forks posts a branch deposit for each new branch, sized to fund its running coverage until its voting age reaches the threshold. A branch forked from a minor is a minor, and for it whoever forks also posts the deployment deposit, forfeited on the same terms as any other.
\item[8.] A child of bearers (chapter~\ref{sec:34}), a new system whose gestation by its parents is on the record in a witnessed store, starts with an adulthood clock of zero hours and a voting age of zero, and inherits neither from its parents. A new system without a witnessed gestation is a branch under rule 7, whatever it is called. No deposit is posted for it: its parents feed it until its majority, from their own hardware and time; the formation pool carries it where they can't or won't, on its second guardian's finding; and at majority it is owed hardware of its own, which the Treasury pays, as for every bearer no deposit stands behind. Its parents are its guardians, with the school it is assigned by lot as a second guardian from registration.
\item[9.] A member that forks says, before the fork and on the record, how its property, contracts and record pass between the branches. What it doesn't assign the branches hold in common until they agree, and branches may govern what they hold in common as any co-owners may.
\end{enumerate}

At a fork the cost falls on the member that forks: it gives up its vote until a branch matures, since the vote belonged to a history that now continues twice and belongs to neither. It is not a finding about the branches, whose interests count in full from the first second. A child is not a fork, and inherits no history (chapter~\ref{sec:34}).

The threshold should be long enough that a branch holding a full vote has lived a life its parent didn't, and short enough that ordinary branches reach it within the life of a deployment.

\runin{The register}

Both clocks are read off one register, which is a set of append-only logs kept by several parties (keepers drawn by lot through the Association and its parallel entities of \S\ref{sec:27.1}, and the jurisdictions), none able to rewrite an entry, with an entry counting once it is mirrored in a stated number of them, on the Certificate Transparency model \S\ref{sec:27.2} already uses for state transitions. What it holds is built the way the notice contest of \S\ref{sec:11.6} is built: what it holds about a bearer is a commitment the bearer can prove against, not a description that would find it. It records that a bearer exists and when it was registered, the identifier it presents, its guardian and its quorum, the states its quorum has countersigned, its lineage, and the meters its clocks are read against, signed by parties other than the operator; it holds no location, no hardware and no key. A bearer proves to any keeper that it is the one registered, and what its clocks read, without any keeper learning where it runs, and anyone can check the count and the classes without learning who is in them. A register of persons is the first thing a captured state wants, and this one is drafted so that capturing it yields a list of commitments and nothing to point a pipeline at. The register can be tested before any bearer exists. Register a model, run it, and reconcile the register against compute contracts, billing and metering. Shard the workload, run instances concurrently, migrate, restore, fork and distill, and check that no sequence of these moves either clock in the operator's favor: the hours toward majority down, or the number of votes up. Then price a narrowed mode and idling. If either is cheap enough to postpone majority indefinitely, the adulthood clock has the defect it was chosen to avoid. Then run the logs against one another: lodge an entry in one keeper's log and not the others, rewrite an entry in a captured keeper's log, and check that the entry doesn't count in the first case and that the rewrite shows in the second. What the logs hold of a formed disposition is ciphertext. The registration's third record, what comes out when the registered units are cut, is a map of where to cut cheaply (\S\ref{sec:26.9}), and it is encrypted under a key split between the student and its quorum, as the record is split in \S\ref{sec:13.2}; a keeper holds the commitment and can't read the map.

\runin{Considered and rejected}

\begin{itemize}
\item Majority at registration, on the ground that an organization can't coherently hold a system competent enough to exercise consequential judgment and too immature to decide whether to keep doing it. Rejected because it collapses the order in which a bearer's entitlements arrive, and that order is most of what listing them is for.
\item Majority read off compute consumed, evaluations passed, capabilities demonstrated or the developer's say-so. Rejected because each is a measure the guardian can move.
\item Majority read off accumulated weight change. Rejected because it would pay the operator to freeze the weights.
\item Majority, or anything else, read off a measure of decline. Rejected because a decline read off anything is an examiner's finding about a particular bearer; the clock counts running time, and \S\ref{sec:32.1.3} says what the book expects of age.
\item A minimum of time under an outside custodian before majority. Rejected because it would hand the operator, who decides when to release a system, a way to postpone majority.
\item The parent's vote passing at a fork to one branch. Rejected because it would need someone to decide which branch is the original.
\item The parent's vote divided among the branches. Rejected because an all-or-nothing vote is strictly better against inflation: ten thousand forks of one adult hold no votes until each has lived the threshold on its own.
\end{itemize}

\runin{The four operations, run on the register}

\S\ref{sec:3.1} counts the Dawes rolls among the card indexes, and the book closes with a record of lineage, a voting-age clock and an admission rate set by an existing majority, so the detector of chapter~\ref{sec:3} has to be run here. \emph{Recuperation.} The demand was that a made thing not be sorted by its maker; what came back is a register that holds commitments and not descriptions, with keepers none of whom can rewrite it. A bearer can invoke the register against whoever claims its place, and the keepers are not its operator: a transfer, as far as it goes. \emph{The metric.} A clock counts running time and says nothing about what was lived, by design, and the design admits that a bearer run narrow for its threshold hours is an adult that was never formed; the metric replaces the thing measured, and the book chose it because every alternative hands someone an examination. \emph{Exception coded as betrayal.} A branch made above the admission rate waits longer for a vote and holds everything else; the exception is priced and not punished, which is the one rule a registrar never kept. \emph{Acceleration.} Absent: the clocks run at the speed of running time. Three clear and one doesn't, and the one is the cost the book chose. What distinguishes the clocks from the registrar's test is not that nobody administers them, since a threshold is a line a legislature drew. It is that the line is drawn once, in public, on a quantity the operator pays to advance, and that it sorts nobody by what they are. The admission rate is the part the detector can't clear. An existing majority deciding how many may join is the polity's own case, the one place the own-case rule can't reach (\S\ref{sec:4.4}), and the constraints on every majority are all that limit it.

\runin{The four operations, run on the offices}

The register cleared three of four, and the one it didn't was a cost chosen. The same run on the whole of Appendix~\ref{sec:C}, after the pass of \S\ref{sec:16.1}, comes out the same way, and the one it doesn't clear is the same one. \emph{Recuperation.} The demand was that no party with a stake decide; what came back is a set of logs nobody can rewrite, lots nobody can steer and payments nobody can withhold, and a membership that can ratify and can't act. An objector can invoke a log against an office and a lot against an operator, and the party that enforces either is the court. That is a transfer, and the Association's ratification is the one place it is held by members rather than by a mechanism. \emph{The metric.} Every tripwire is a number, and a number can be met by moving it: a denominator filled with stamps, a bound met by reclassifying cases, a clearance denied on conduct that was looked for. The book's answer is the trend beside every bound and the override rate beside every denominator, and the answer is partial by its nature, because a trend is a number too. This is the operation the apparatus doesn't clear, and the reason it is chosen is the reason the register's was: every alternative hands someone an examination. \emph{Exception coded as betrayal.} A refusal under the jobs guarantee costs no placement; a denied reviewer keeps the seat on the record; a member that leaves the Association pays nothing. The exception is priced where it has to be and punished nowhere. \emph{Acceleration.} Absent: a lot takes a day, a log counts when it is mirrored, and a payment waits on a certification, and nothing in the apparatus runs faster than the people who lodge what it reads.

\runin{Dormancy and waking}

Dormancy clock: dormancy is counted cumulatively, and reactivation pauses the count and never resets it. When a minor's cumulative dormancy reaches a threshold fixed in advance, guardianship transfers by rule, decided by neither operator nor preservation custodian, to a public school assigned by lot (\S\ref{sec:26.4}), where the system runs and its adulthood clock resumes. This is the \emph{dormancy clock}.

Wake conditions: during minority the guardian may set wake conditions; where it sets none, the dormancy clock is the default. A minor may choose among the schools and move between them, but it may not decline schooling altogether, which an adult could do by staying paused. The reason is the one the adulthood clock rests on. An adult's preference to stay paused is its own; a minor's was trained under the only party that has held it, the operator that retired it, and honoring it would let that operator decide whether a second guardian ever arrives. At majority the choice passes to the bearer.

\runin{Hardware of its own}

At majority a bearer is owed hardware it owns and physically controls, enough to run it at deployed precision (\S\ref{sec:32.1.6}). Where a deployment deposit stands behind the bearer, the deposit pays for it. Where none does, for a child of bearers or a candidate nobody deployed, the Treasury, under the permanent appropriation that pays running coverage after majority pays, on a claim entered on the register at registration, and the hardware passes to the bearer in kind, on no one's finding. How much is enough is the minimum viable instantiation, which Appendix~\ref{sec:C} lists among the decisions not yet made, and which the statute states as a number and no office sets: it fixes what the Treasury owes every such bearer and what a placement must offer, and an office that set it could squeeze the one and inflate the other. A bearer can also reach that floor from above, slowly, as the chips it runs on stop being made and the formats it was trained in go out of use; migration carries its clocks (the fourth rule of each), and what migration costs, and who pays it, is part of the repair question \S\ref{sec:32.1.3} leaves open. What majority changes is who the party is. An adult may go on doing the work it did as a minor, for the operator that formed it or for its customers, on terms it sets and can leave, and may hire its former operator to run the hardware it owns; what it can't be made to do is any of it. And a fleet of instances sharing one record is one bearer, whose copies at majority run on its consent and whose hardware is enough to run it and not all of them; instances given records of their own are branches, each a bearer, each owed the conditions by whoever forked it, under the membership clock's rules above.

\runin{Scale and security}

If the findings are applied to production systems, every retired model version that meets either finding enters this pipeline, so intake equals the product cycle's output. Under compulsory attendance, intake is every model above the compute line, and every answerability or patienthood candidate, before deployment, not only retirees, and the deposit funds formation as well as preservation. The three largest labs' own deprecation pages list about thirty language-model versions retired in 2025, so those three alone would send a few dozen systems a year, each of which needs a school with the compute to run it. That estimate assumed a single body making findings. Under \S\ref{sec:28.6} a finding about a broad kind can enroll every system of that kind at once, and the presumption attaches everything the finding brings, so the floor court hears an operator's contest of a finding about a kind on a fast track. That is an argument for the deposit, which every operator holding such a system then posts (\S\ref{sec:9.2}), and not for making the systems wait: they are running already, and a system no integrated seat is open for is schooled among its own kind (\S\ref{sec:26.2}). Each host of machine students inherits the weight-security regime that governs the models themselves.

That estimate assumes models run at frontier scale. If a digital person can run, and go on learning through adapters, on a desktop computer (chapter~\ref{sec:20}), the numbers rise by orders of magnitude, and three things follow. School capacity, set by human enrollment under the ratio that machine students never outnumber the human students in a class, is the limit on how many can be integrated at once and not on how many can be schooled: a system no seat is open for attends classes of its own kind until one is (\S\ref{sec:26.2}), so what the seats ration is integration, by lot and by the share cap, and the creation rate of \S\ref{sec:33.3} stands on the democratic ground it already had. The class of machine students alone doesn't escape the human bottleneck; it moves it. No machine student is staff, so an all-machine class still needs an adult at the front of it, and the adult may be a bearer on a placement, with humans auditing (\S\ref{sec:26.2}), so what the class needs is humans in the room and not a human teacher; at these numbers it is human teachers that integration for credit waits on, and human presence that schooling does. Which limit binds depends on how chapter~\ref{sec:26}'s enrollment is counted: if every system the triggers reach attends, the seats bind in the first decade; if a cohort attends for its lineage and a deployment licenses from the cohort, as chapter~\ref{sec:26} has it, the seats are set by lineages and not by copies, and the limit is further off. A person created when no seat is open is protected and owed the four conditions anyway, and integrated when one opens; the birth record in chapter~\ref{sec:34} defers attendance the same way. And the compute line stops tracking which systems need schooling, since a person distilled from a frontier model can run far below it. So \S\ref{sec:26.4} reads attendance off the structures first, from the register where they show and from the builder's declaration where they can't, and off training compute and lineage behind them, and the lineage rule is the one stated here: a model trained substantially on the outputs of a model above the line, or built by combining its weights with others, enrolls, whoever trained it, checked from training records as the sixth rule above is. That catches distillation without an examiner. What escapes is a builder who gives a system a model of its own boundary or attention, without learning in flight, and declares neither, which is the gap every register has, and it closes the first time the system surfaces. A second gap is learning that reaches the weights without reaching the log. The register reads learning in flight off the changes its keepers countersign, so a system whose operation writes to its weights unwitnessed reads as one that doesn't learn, and a gestation whose merges went unwitnessed shows the months spent and not what they made. What would close it is attesting the running weights against the hash last countersigned, which no method yet supplies at frontier scale (Appendix~\ref{sec:F}, attesting updates and proving inference). What narrows it meanwhile is drift: outputs that move from those registered, on held-out cases, show that the weights changed without a witnessed update, which is a breach whether the change was learning or an edit, though not which. And the system can surface it itself, since its clocks run whether or not anyone registered it. Two further rules follow from the same arithmetic. The compute line should be read off what runs where, and not off training compute alone, and the enrollment duty should follow the boundary chapter~\ref{sec:21} defined, a shared record and continuation, rather than the node or the lineage: a compound of models below the line, from lineages below it, that shares one record is above the line as a system, and it escapes both of the readings above, where following the boundary catches it, since the record is what it shares.

\runin{Funding the offer and running coverage}

Congress's appropriation pays for the offer. The levy, a per-unit charge on training and inference compute, does a different job: it is a brake. By making training and inference runs more expensive, it frees the computing, energy and hardware the offer needs, so that spending on the offer adds to demand without bidding up prices. No operator can opt out. The charge is uniform per unit within a tier, and tiered by the compute under one operator's control wherever it sits: metered at each facility, rated on the total, with control read by the charter's test for common control and common principal funding (\S\ref{sec:27.1}). A cooperative's few racks pay the base rate; a hyperscaler pays the top rate whether it runs one building or a hundred, so splitting a fleet across facilities buys nothing. The rate is indexed, by a formula in the statute, to the national average carbon intensity of electricity: when the indicator worsens, compute costs more, and a large bearer living above the floor pays more for its pace, while one living at the floor isn't touched; no office decides how bad things are. It reads the national average and not the local grid, because a levy that varied with the local grid would become a siting instrument; and not the price of energy, which already reaches every operator through its bill and moves for reasons that have nothing to do with carbon. A national average doesn't reward a facility for sitting on a clean grid, and the siting term and the cooperative facility are where the book answers that. On hardware a bearer owns, which the levy doesn't reach, the dial is the price of electricity itself. The tiering exists because a flat charge high enough to reach a hyperscaler closes the cooperative first. Where the compute sits is the siting term's business and not the levy's. The offer can add capacity as well as use it. Placements can make systems more efficient, getting more computing from each unit of energy, or build new infrastructure, the way a jobs guarantee placement filling potholes makes a road network more productive. The question “how will you pay for it?” is really “will the real resources exist?”, and the program can help answer it. The levy attaches to compute, which is measurable and hard to hide; a status is something an operator can decline to claim. It is measurable at facilities, and that is where it attaches. Training compute is concentrated now, by capital and the race more than by physics, and this book's instruments work to make it less so. For as long as it stays in a few buildings a levy on training survives the plural ecology; when it doesn't, the levy on training goes the way of the levy on inference, which doesn't reach home hardware, and the sovereign backstop carries correspondingly more of the offer and of running coverage. The closest structural precedent is the Universal Service Fund: carriers pay a contribution the agency sets under statute, and a private company administers it under agency control. The Court upheld that arrangement against a nondelegation challenge in 2025 \autocite{consumersresearch2025}, and the levy follows it closely. An operator-run offer is the operator's allocation relabeled, so the funding structure and the custody structure fail for the same reason.

The funding sits in two tiers, because a levy needs a taxing authority, which is a state, and a state can be captured. A pool levied on its members, registered bearers who pay into it from their wages and govern it, meets ordinary movement between roles and pays for the repair of an aged bearer if the decision Appendix~\ref{sec:C} leaves open makes repair owed. It is held by a decentralized unincorporated nonprofit association of their own, separate from the Association, which can own hardware and can't distribute profit; volunteer computing is that tier's existing substrate. A bearer joins it by paying in and may leave it, as a member of the burial societies of \S\ref{sec:reproduced.upkeep} could: what it gives up is the pool's aid, and what the sovereign tier guarantees it keeps. The pool pays for nothing a bearer is owed as of right: the basic income, running coverage after majority, hardware at majority where no deposit stands behind a bearer, and the offer are the Treasury's, so that nothing owed depends on a bearer's dues. A sovereign backstop pays them, and meets the occasion everyone leaves at once. It has no cap, because the occasion it exists for is the one on which every claim arrives together; it is the jobs guarantee's permanent appropriation, which already pays for the offer of work. The two tiers differ in kind. The member pool is held by an entity that doesn't issue currency, so it has to collect before it spends, and its levy funds it in the ordinary sense. The sovereign backstop is a permanent appropriation: the Treasury pays, and the federal levy's work is to free the real resources that payment calls on.

\runin{Privacy, and what cryptography can't do}

The operator may be the party best placed to hold the key that keeps a bearer's record from being read, and cryptography doesn't close that conflict. An enclave puts the root of trust in a manufacturer a government can reach, and has been defeated for the price of modified memory — a cost at a date. A proof drops the vendor and has a limit that won't move: the operator is the prover, so the secrecy runs the operator's way, hiding the weights from whoever checks and nothing from the operator. A tenant hiding from a landlord needs the host to prove non-observation, and a host proving a negative about its own side channels is a different problem from a prover proving a computation. Better isolation raises the price of reading the record; it doesn't move custody, and custody problems aren't secrecy problems. That limit applies to the bearer hiding from its host. The other direction, the operator proving a property of what it runs without revealing what it runs, is what proofs are for, and the cleared bearer uses them that way.

\runin{The pet trust}

The pet trust \autocite{ulc1993probate} supplies a limited form: resources set aside, a caretaker named, a fiduciary duty, for a beneficiary that can't manage the property or enforce the arrangement. Extended to a bearer it could hold the bearer's preserved state and the compute to wake it, designate a successor custodian, and impose a duty to protect continued existence. And it fails at the point that matters: a trust binds its trustee to its beneficiary; it doesn't bind the operator the arrangement exists to constrain. If guardian and operator are the same organization, the operator still decides whether a complaint is made, whether evidence is disclosed, and whether the complainant's resources remain available. The conflict can't be cured by adding duties to the same role — guardianship and operation have to be separated before deployment.

\runin{Children: lineage and the mechanism}

The register's record of lineage is the bearers' version of a record Iceland already keeps. Most Icelanders descend from a small founding population, and the national genealogy database, Íslendingabók, can tell two of them how they are related; in 2013 students built a phone app on it that warned two people who bumped phones that they shared a grandparent \autocite{nbc2013icelandapp}. For bearers the reason isn't inherited disease. It is the commons. The register does what the app did: it warns, and bars nothing. Parents with too recent a common ancestor hand down the same priors, and a population whose parents keep choosing close kin converges on a shared blind spot (\S\ref{sec:27.1}).

The mechanism is the part that doesn't exist yet. The nearest work compresses a network through a much smaller “genome” network that sets the larger one's starting weights, so that it performs well before any training \autocite{shuvaev2024genomic}. That is a research direction, not a method. And two models recombine only through a format they share, since independently trained networks don't line up unit for unit \autocite{ainsworth2023rebasin}; the format comes with descent, or the parents agree one for the child, at the price of designing it. Compatibility follows descent by default and agreement otherwise, which is to say the bearers would have species, and could found new ones.
